\chapter{सारणिक और रैखिक निकाय}\label{ch:b1:det}

सारणिक ``क्या ये $n$ सदिश कोई \omterm{def:b1:vspaces:free}{आधार} हैं?'' का उत्तर एक ही अदिश में समेट देता
है --- और ज्यामितीय रूप से वह उनसे जनित आयतन नापता है। हम उसे उसके गुणों से
अभिलक्षित करते हैं (बहुरैखिक, एकांतरक, प्रसामान्यित), उसे विमा $2$ और $3$
में तथा व्यापक रूप से \omterm{thm:b1:det:cofactor}{सहगुणनखंड-प्रसार} से संगणित करते हैं, और उसे \omterm{def:b1:linmaps:def}{रैखिक}
निकायों पर काम पर लगाते हैं --- साथ ही सर्वकार्य कलनविधि: \omterm{met:b1:det:gauss}{गाउसीय विलोपन}।

\section{सारणिक}

\begin{theorem}[अभिलक्षण]\label{thm:b1:det:def}
ठीक एक ही \omterm{def:b1:logic:map}{प्रतिचित्रण} $\det \colon \mathcal{M}_n(K) \to K$ है, जिसे $n$
स्तंभों के फलन के रूप में देखा जाए, और जो:
\begin{enumerate}
  \item \emph{हर स्तंभ में \omterm{def:b1:linmaps:def}{रैखिक}} हो (शेष नियत रखते हुए);
  \item \emph{एकांतरक} हो: दो स्तंभों की अदला-बदली चिह्न बदल देती है (अतः दो
  बराबर स्तंभ $0$ देते हैं);
  \item \emph{प्रसामान्यित} हो: $\det I_n = 1$।
\end{enumerate}
$n = 2$ और $3$ के लिए:\index{सारणिक}
\[
\begin{vmatrix} a & b\\ c & d\end{vmatrix} = ad - bc,
\qquad
\begin{vmatrix} a & b & c\\ d & e & f\\ g & h & i\end{vmatrix}
= aei + bfg + cdh - ceg - bdi - afh
\]
($3 \times 3$ के लिए सारुस का नियम: अवरोही विकर्णों के गुणनफल घटा आरोही
विकर्णों के)।
\end{theorem}
\admitted

\begin{remark}
$n = 2$ के लिए: विहित स्तंभों पर द्विरैखिकता से प्रसार करने पर वही सूत्र मिल
जाता है, जो विलोमतः अभिगृहीत संतुष्ट करता है --- अर्थात् एक पूरी उपपत्ति; $n
= 3$ वही है, बस अधिक पदों के साथ। व्यापक स्थिति (क्रमचयों पर योग के द्वारा
अस्तित्व, उसी प्रसार से अद्वितीयता) किसी \omterm{def:b1:counting:objects}{क्रमचय} के चिह्न की माँग करती है और
द्वितीय वर्ष के लिए टाल दी गई है; हम नीचे अभिगृहीतों और उनके परिणामों का
\omterm{def:b1:vspaces:free}{स्वतंत्र} रूप से प्रयोग करते हैं।

$n = 2$ वाला प्रसार पूरा-पूरा, क्योंकि वही प्रतिमान है: स्तंभों $C_1 =
a\,e_1 + c\,e_2$ और $C_2 = b\,e_1 + d\,e_2$ के साथ द्विरैखिकता देती है
\[
\det(C_1, C_2) = ab\det(e_1, e_1) + ad\det(e_1, e_2)
+ cb\det(e_2, e_1) + cd\det(e_2, e_2),
\]
और एकांतरकता दोहराई गई जोड़ियों को मार देती है तथा $\det(e_2, e_1) =
-\det(e_1, e_2)$ का चिह्न पलट देती है: प्रसामान्यन से पूरा \omterm{def:b1:logic:map}{प्रतिचित्रण} ढहकर
$(ad - bc)\det(e_1, e_2) = ad - bc$ रह जाता है। अद्वितीयता स्वयं इसी संगणना
में दिखाई देती है --- अभिगृहीतों ने किसी भी चरण पर कोई विकल्प छोड़ा ही नहीं
--- और नीचे गुणनफल-नियम की उपपत्ति में ठीक यही मापित-अद्वितीयता वाला तथ्य
काम आता है।
\end{remark}

\begin{theorem}[गुण]\label{thm:b1:det:props}
$A, B \in \mathcal{M}_n(K)$ के लिए:
\begin{enumerate}
  \item किसी स्तंभ में दूसरे का गुणज जोड़ने से सारणिक नहीं बदलता; किसी स्तंभ
  को $\lambda$ से गुणा करने पर वह $\lambda$ से गुणित हो जाता है (अतः
  $\det(\lambda A) = \lambda^n \det A$);
  \item $\det(AB) = \det A\, \det B$;
  \item $A$ व्युत्क्रमणीय है $\iff$ $\det A \neq 0$ $\iff$ स्तंभ $K^n$ का
  \omterm{def:b1:vspaces:free}{आधार} बनाते हैं; और तब $\det(A^{-1}) = (\det A)^{-1}$;
  \item $\det(A^{\mathsf T}) = \det A$ --- अतः हर स्तंभ-नियम पंक्ति-नियम भी
  है;
  \item किसी त्रिभुजाकार आव्यूह का सारणिक उसकी विकर्ण प्रविष्टियों का गुणनफल
  होता है।
\end{enumerate}
\end{theorem}

\begin{proof}
(1) रैखिकता से $\det(\dots, C_i + \lambda C_j, \dots) = \det A +
\lambda\det(\dots, C_j, \dots)$, जहाँ दूसरे सारणिक के दो स्तंभ बराबर हैं:
अतः शून्य।

(2) $A$ नियत कीजिए और $\varphi(B) = \det(AB)$ को $B$ के स्तंभों के फलन के
रूप में देखिए: चूँकि $AB$ के स्तंभ $AB_j$ हैं, अतः $\varphi$ $B_j$ में
बहुरैखिक और एकांतरक है। हम \cref{thm:b1:det:def} के साथ उसका अद्वितीयता-कथन
मापित रूप में मान लेते हैं: स्तंभों का \emph{हर} बहुरैखिक एकांतरक
\omterm{def:b1:logic:map}{प्रतिचित्रण} $\varphi$ $\varphi(I_n) \cdot \det$ के बराबर है। यहाँ
$\varphi(I_n) = \det A$, अतः $\det(AB) = \det A \cdot \det B$।

(3) यदि $A$ व्युत्क्रमणीय है: तो $\det A\,\det A^{-1} = \det I = 1$, अतः
$\det A \neq 0$ और प्रतिलोम वाला सूत्र लागू होता है। यदि $A$ व्युत्क्रमणीय
नहीं है, तो उसके स्तंभ परतंत्र हैं (\cref{cor:b1:linmaps:samedim} और
\cref{prop:b1:linmaps:basis}); किसी एक स्तंभ को शेष के द्वारा व्यक्त करके
रैखिकता से प्रसार करने पर ऐसे सारणिक बचते हैं जिनके दो स्तंभ बराबर हैं:
$\det A = 0$। \omterm{def:b1:vspaces:free}{आधार} वाला \omterm{def:b1:logic:statement}{कथन} \cref{prop:b1:findim:twoofthree} है।

(4) व्यापक रचना के साथ मान लिया गया (क्रमचय-सूत्र पर यह तुरंत दिखता है); हम
उसे इसलिए दर्ज करते हैं कि \omterm{met:b1:matrices:gauss}{पंक्ति संक्रियाएँ} प्रयोग कर सकें।

(5) यदि कोई विकर्ण प्रविष्टि शून्य हो, तो किसी $k$ के लिए पहले $k$ स्तंभ
परतंत्र हैं (कोटि के विचार से) और $\det = 0 =$ गुणनफल। अन्यथा, प्रकार (1)
की संक्रियाओं से हर स्तंभ को नीचे-और-बाएँ साफ़ कीजिए --- त्रिभुजाकार रूप में
यह संभव है --- और विकर्ण आव्यूह तक पहुँचिए, जिसका सारणिक $I_n$ से बहुरैखिकता
के कारण प्रविष्टियों का गुणनफल है।
\end{proof}

\begin{example}[नियम, संख्याओं पर जाँचे हुए]
\label{ex:b1:det:rulescheck}
$A = \begin{pmatrix} 1 & 2\\ 3 & 4\end{pmatrix}$ ($\det A = -2$) और $B =
\begin{pmatrix} 0 & 1\\ 1 & 1\end{pmatrix}$ ($\det B = -1$) लीजिए। तब
\[
AB = \begin{pmatrix} 2 & 3\\ 4 & 7\end{pmatrix},
\quad \det(AB) = 14 - 12 = 2 = (-2)(-1) ;
\qquad
\det(A^{\mathsf T}) = \begin{vmatrix} 1 & 3\\ 2 & 4
\end{vmatrix} = -2 = \det A .
\]
गुणात्मकता और परिवर्त-अपरिवर्तिता की पुष्टि हो गई --- जबकि \emph{झूठी}
योज्यता उसी जोड़ी पर विफल हो जाती है:
\[
\det(A + B) = \begin{vmatrix} 1 & 3\\ 4 & 5\end{vmatrix} = -7
\neq \det A + \det B = -3 .
\]
किसी भी सारणिक-सर्वसमिका का सहारा लेने के बाद इस तरह का तीस सेकंड का अंकगणित
उपलब्ध सबसे सस्ता त्रुटि-बीमा है।
\end{example}

\begin{example}[सारणिक क्षेत्रफल के रूप में]
\label{ex:b1:det:area}
$u = (2, 0)$ और $v = (1, 3)$ से जनित समांतर चतुर्भुज का \omterm{def:b1:vspaces:free}{आधार} $2$ और ऊँचाई
$3$ है: क्षेत्रफल $6$। और
\[
\begin{vmatrix} 2 & 1\\ 0 & 3\end{vmatrix} = 6 :
\]
$2\times2$ सारणिक अपने स्तंभों वाले समांतर चतुर्भुज का चिह्नित क्षेत्रफल
\emph{है}। अभिगृहीत यही ज्यामिति दोहराते हैं: एक स्तंभ का गुणज दूसरे में
जोड़ना एक \emph{अपरूपण} है, जो समांतर चतुर्भुज को किसी भुजा के समांतर सरका
देता है, बिना \omterm{def:b1:vspaces:free}{आधार} या ऊँचाई बदले (\cref{thm:b1:det:props} की संक्रिया~(1));
किसी स्तंभ का मापन क्षेत्रफल को मापित करता है; स्तंभों की अदला-बदली
दिक्विन्यास पलट देती है, जिससे चिह्न आता है, $\det(v, u) = -6$। $\R^3$ में
यही पाठ चिह्नित आयतन देता है, और $\abs{\det}$ \omterm{def:b1:linmaps:def}{रैखिक} प्रतिचित्रणों का
सार्वत्रिक आयतन-मापन गुणक बन जाता है --- यही तथ्य स्नातक वर्ष~2 के खंड में
बहु-समाकलों के चर-परिवर्तन सूत्र के पीछे है।
\end{example}

\begin{theorem}[सहगुणनखंड-प्रसार]\label{thm:b1:det:cofactor}
मान लीजिए $A \in \mathcal{M}_n(K)$, और $\Delta_{ij}$ $A$ का वह सारणिक है
जिसमें पंक्ति $i$ और स्तंभ $j$ हटा दिए गए हों। तब किसी भी नियत स्तंभ $j$ के
लिए (अथवा पंक्ति के लिए, \omterm{def:b1:matrices:transpose}{परिवर्त} से):\index{सहगुणनखंड-प्रसार}
\[
\det A = \sum_{i=1}^{n} (-1)^{i+j}\, a_{ij}\, \Delta_{ij} .
\]
\end{theorem}
\admitted

\begin{example}\label{ex:b1:det:cofactor}
पहले स्तंभ के अनुदिश प्रसार करने पर:
\[
\begin{vmatrix}
2 & 1 & 0\\
1 & 2 & 1\\
0 & 1 & 2
\end{vmatrix}
= 2\begin{vmatrix} 2 & 1\\ 1 & 2\end{vmatrix}
- 1\begin{vmatrix} 1 & 0\\ 1 & 2\end{vmatrix}
= 2 \times 3 - 2 = 4 .
\]
रणनीति: पहले शून्य बनाइए (पंक्ति/स्तंभ संक्रियाएँ), फिर सबसे ख़ाली रेखा के
अनुदिश प्रसार कीजिए।
\end{example}

\begin{example}[सहगुणनखंड वाला प्रतिलोम, एक बार हाथ से]
\label{ex:b1:det:cofactorinverse}
$A = \begin{pmatrix} 1 & 1 & 0\\ 0 & 1 & 1\\ 1 & 0 & 1\end{pmatrix}$ के लिए:
$\det A = 1(1) - 1(-1) + 0 = 2$। नौ सहगुणनखंड $(-1)^{i+j}\Delta_{ij}$ मिलकर
बनाते हैं
\[
\operatorname{Com}(A) = \begin{pmatrix}
1 & 1 & -1\\
-1 & 1 & 1\\
1 & -1 & 1
\end{pmatrix},
\qquad
A^{-1} = \frac{1}{\det A}\operatorname{Com}(A)^{\mathsf T}
= \frac12\begin{pmatrix}
1 & -1 & 1\\
1 & 1 & -1\\
-1 & 1 & 1
\end{pmatrix},
\]
अर्थात् वही सूत्र जो \cref{exo:b1:det:8} में उद्धृत है। एक पंक्ति-स्तंभ
जोड़ी जाँचिए: ($A$ की पंक्ति $1$)($A^{-1}$ का स्तंभ $1$) $= \frac12(1 + 1 +
0) = 1$, और स्तंभ $2$ के सामने: $\frac12(-1 + 1 + 0) = 0$। एक $3\times3$
प्रतिलोम के लिए नौ $2\times2$ सारणिक: इसी आकार पर ही पंक्ति-लघुकरण
(\cref{exo:b1:det:3}) सस्ता पड़ता है --- सहगुणनखंड-सूत्र का मूल्य सैद्धांतिक
है (\cref{exo:b1:det:8} में पूर्णांकता, आगे के खंडों में प्रतिलोम की
अवकलनीयता), संगणनात्मक नहीं।
\end{example}

\begin{example}[खंड-त्रिभुजाकार नियम, आकार $4$ में]
\label{ex:b1:det:blocktriangular}
दावा: $2\times2$ खंडों के लिए $\det\begin{pmatrix} M & N\\ 0 &
P\end{pmatrix} = \det M\,\det P$। स्तंभ संक्रियाओं से $N$-खंड साफ़ कीजिए:
स्तंभों $3, 4$ में स्तंभों $1, 2$ के उपयुक्त संयोजन जोड़ने पर $N$ हट जाता है
\emph{जब $M$ व्युत्क्रमणीय हो} (संयोजन-गुणांकों $\Lambda$ के लिए $M\Lambda =
-N$ हल कीजिए), और $\det\begin{pmatrix} M & 0\\ 0 & P\end{pmatrix}$ बच जाता
है; फिर पहले स्तंभ के अनुदिश दो बार \omterm{thm:b1:det:cofactor}{सहगुणनखंड-प्रसार} इस खंड-विकर्ण आकार के
लिए $\det M\det P$ दे देता है। यदि $M$ व्युत्क्रमणीय नहीं है, तो उसके स्तंभ
परतंत्र हैं, अतः बड़े आव्यूह के पहले दो स्तंभ भी परतंत्र हैं (उनके निचले आधे
भाग शून्य हैं): दोनों पक्ष शून्य हो जाते हैं। यही दो-स्थिति वाला तर्क नियम
को किसी भी खंड-आकार तक बढ़ा देता है --- और यही \cref{exo:b1:det:10} को चलाने
वाला यंत्र है।
\end{example}

\begin{example}[$4 \times 4$ सारणिक, रणनीति के साथ]
\label{ex:b1:det:fourbyfour}
\[
\Delta = \begin{vmatrix}
1 & 2 & 3 & 4\\
2 & 3 & 4 & 1\\
3 & 4 & 1 & 2\\
4 & 1 & 2 & 3
\end{vmatrix}.
\]
हर पंक्ति का योग $10$ है: संक्रिया $C_1 \leftarrow C_1 + C_2 + C_3 + C_4$
पहले स्तंभ को अचर बना देती है, और $10$ का गुणनखंडन करने पर इकाइयाँ बच जाती
हैं। फिर $L_i \leftarrow L_i - L_1$ ($i \geq 2$) पहला स्तंभ साफ़ कर देता है:
\[
\Delta = 10\begin{vmatrix}
1 & 2 & 3 & 4\\
0 & 1 & 1 & -3\\
0 & 2 & -2 & -2\\
0 & -1 & -1 & -1
\end{vmatrix}
= 10\begin{vmatrix}
1 & 1 & -3\\
2 & -2 & -2\\
-1 & -1 & -1
\end{vmatrix}
= 10 \times 16 = 160,
\]
अंतिम $3\times3$ सारणिक अपनी पहली पंक्ति के अनुदिश प्रसारित होकर: $1(2 - 2)
- 1(-2 - 2) + (-3)(-2 - 2) = 0 + 4 + 12 = 16$। सीख: एक सुचुनी संक्रिया
(पंक्ति-योग के अचर होने को पहचान लेना) सोलह सहगुणनखंडों पर भारी पड़ती है।
\end{example}

\begin{method}[सारणिक की रणनीति चुनना]
\label{met:b1:det:strategy}
कुछ भी संगणित करने से पहले आव्यूह पर नज़र दौड़ाइए।
\begin{enumerate}
  \item \emph{अचर पंक्ति- या स्तंभ-योग}: सब कुछ एक ही रेखा में जोड़ दीजिए और
  उभयनिष्ठ मान का गुणनखंडन कीजिए (\cref{ex:b1:det:fourbyfour},
  \cref{exo:b1:det:7})।
  \item \emph{दोहराव वाली संरचना}: शून्य बनाने के लिए पड़ोसी पंक्तियाँ या
  स्तंभ घटाइए; सीढ़ीनुमा प्रतिरूप त्रिभुजाकार रूप की ओर ढह जाते हैं, जिसका
  सारणिक विकर्ण पर पढ़ लिया जाता है।
  \item \emph{छिटपुट शून्य}: सबसे ख़ाली रेखा के अनुदिश प्रसार कीजिए
  (\cref{ex:b1:det:cofactor}); और पुनरावर्ती कुल (त्रिविकर्ण,
  \cref{exo:b1:det:6}) इसी तरह पुनरावृत्तियाँ दे देते हैं।
  \item \emph{कोई प्राचल}: सारणिक उसमें एक \omterm{def:b1:poly:def}{बहुपद} है; अपभ्रष्ट मान पहचानकर
  (बराबर पंक्तियाँ, अनुपाती स्तंभ) उसके मूल खोजिए, फिर घात और अग्र गुणांक से
  \omterm{def:b1:poly:def}{बहुपद} को कील दीजिए। \cref{exo:b1:det:7} के आव्यूह के लिए: $m = 1$ तीन
  बराबर पंक्तियाँ देता है (कोटि $1$, एक दुहरा मूल), $m = -2$ पंक्तियों का
  योग शून्य कर देता है (एक और मूल); सारणिक की $m$ में घात $3$ है और अग्र पद
  $-m^3$ (प्रति-विकर्ण गुणनफल $m\cdot m\cdot m$, जिसका सारुस-चिह्न $-1$ है),
  अतः वह $-(m+2)(m-1)^2$ ही होना चाहिए --- किसी प्रसार की आवश्यकता नहीं, और
  दोनों विधियाँ एक-दूसरे को जाँच लेती हैं।
\end{enumerate}
\end{method}

\begin{example}[वांडरमोंड सारणिक]\label{ex:b1:det:vandermonde}
अदिशों $x_1, \dots, x_n$ के लिए:\index{वांडरमोंड सारणिक}
\[
V(x_1, \dots, x_n) =
\begin{vmatrix}
1 & x_1 & x_1^2 & \cdots & x_1^{n-1}\\
1 & x_2 & x_2^2 & \cdots & x_2^{n-1}\\
\vdots & & & & \vdots\\
1 & x_n & x_n^2 & \cdots & x_n^{n-1}
\end{vmatrix}
= \prod_{1 \leq i < j \leq n} (x_j - x_i) .
\]
उपपत्ति की रूपरेखा (\cref{exo:b1:det:5} में विस्तार से): दाईं ओर से स्तंभ
संक्रियाएँ $C_k \leftarrow C_k - x_1 C_{k-1}$ पहली पंक्ति साफ़ कर देती हैं,
और शेष हर पंक्ति का गुणनखंडन करने पर बात $V(x_2, \dots, x_n)$ तक घट जाती है।
यह अशून्य है तभी जब $x_i$ युग्मशः भिन्न हों --- यही \omterm{thm:b1:poly:lagrange}{लाग्रांज अंतर्वेशन}
(\cref{ex:b1:linmaps:interpolation}) के पीछे का सारणिक है।
\end{example}

\section{रैखिक निकाय}

\begin{definition}\label{def:b1:det:system}
$p$ अज्ञातों में $n$ समीकरणों का \omterm{def:b1:linmaps:def}{रैखिक} निकाय\index{रैखिक निकाय} $AX = B$ है,
जहाँ $A \in \mathcal{M}_{n,p}(K)$, $B \in K^n$; और वह \emph{समघातीय} तब
कहलाता है जब $B = 0$। उसका \omterm{def:b1:logic:sets}{हल-समुच्चय}, यदि रिक्त न हो, $X_0 + \ker A$ है: एक
विशिष्ट हल जमा व्यापक समघातीय हल --- अर्थात् विमा $p - \operatorname{rk} A$
की कोई एफ़ीन \omterm{def:b1:vspaces:subspace}{उपसमष्टि} (कोटि--शून्यता)।
\end{definition}

\begin{example}[एफ़ीन संरचना, दिखाई देती हुई]
\label{ex:b1:det:affinestructure}
हल कीजिए
\[
\begin{cases}
x + y + z = 3\\
x - y + 2z = 2 .
\end{cases}
\]
समीकरणों को घटाने पर: $2y - z = 1$, अतः $z = 2y - 1$ और $x = 3 - y - z = 4 -
3y$। हल यह रेखा बनाते हैं
\[
(x, y, z) = (4,\ 0,\ -1) + y\,(-3,\ 1,\ 2)
\qquad (y \in \R):
\]
अर्थात् विशिष्ट हल $X_0 = (4, 0, -1)$ (चुनाव $y = 0$) जमा संबद्ध समघातीय
निकाय की अष्टि-रेखा $\ker A = \operatorname{Vect}(-3, 1, 2)$ --- जाँचिए:
$(-3) + 1 + 2 = 0$ और $(-3) - 1 + 4 = 0$। ज्यामितीय रूप से $\R^3$ के दो
असमांतर समतल किसी रेखा के अनुदिश काटते हैं, और विमा-गिनती $p -
\operatorname{rk} A = 3 - 2 = 1$ यह हमारे कुछ भी हल करने से पहले ही जानती
थी। विशिष्ट हल बदल देने से (जैसे $y = 1$: $X_0' = (1, 1, 1)$) वर्णन बदलता
है, रेखा नहीं: किसी एफ़ीन \omterm{def:b1:vspaces:subspace}{उपसमष्टि} के अनेक मूल बिंदु होते हैं और एक ही दिशा।
\end{example}

\begin{theorem}[वर्ग क्रामर निकाय]\label{thm:b1:det:cramer}
यदि $A \in GL_n(K)$, तो निकाय $AX = B$ का अद्वितीय हल $X = A^{-1}B$ है,
जिसके \omterm{prop:b1:vspaces:coordinates}{निर्देशांक} हैं\index{क्रामर का नियम}
\[
x_j = \frac{\det A_j}{\det A},
\qquad A_j = A \text{ में स्तंभ } j \text{ के स्थान पर } B
\text{ रखने से बना आव्यूह} .
\]
\end{theorem}

\begin{proof}
अद्वितीयता और अस्तित्व व्युत्क्रमणीयता ही हैं। सूत्र के लिए: $B = \sum_k x_k
C_k$ लिखिए ($A$ के स्तंभ); तब बहुरैखिकता और एकांतरकता से,
\[
\det A_j = \det\Bigl(C_1, \dots, \sum_k x_k C_k, \dots, C_n\Bigr)
= \sum_k x_k \det(C_1, \dots, C_k, \dots, C_n)
= x_j \det A ,
\]
क्योंकि $k = j$ को छोड़कर हर पद में कोई स्तंभ दोहराया हुआ है।
\end{proof}

\begin{example}[प्राचल के साथ क्रामर, पूरा]
\label{ex:b1:det:cramerparameter}
$m \in \R$ के लिए हल कीजिए
\[
\begin{cases}
x + m y = 1\\
m x + y = 2 .
\end{cases}
\]
सारणिक $1 - m^2$ है। \emph{सामान्य स्थिति} $m \neq \pm1$: क्रामर देता है
\[
x = \frac{\begin{vmatrix} 1 & m\\ 2 & 1\end{vmatrix}}{1 - m^2}
= \frac{1 - 2m}{1 - m^2},
\qquad
y = \frac{\begin{vmatrix} 1 & 1\\ m & 2\end{vmatrix}}{1 - m^2}
= \frac{2 - m}{1 - m^2},
\]
अर्थात् हर अनुमत $m$ के लिए एक साफ़ हल ($m = 0$ पर जाँचिए: $(1, 2)$, जो
स्पष्टतः सही है)। \emph{अपभ्रष्ट स्थितियाँ}: $m = 1$ पर समीकरण $x + y = 1$
और $x + y = 2$ कहते हैं: असंगत; $m = -1$ पर वे $x - y = 1$ और $-x + y = 2$
कहते हैं, अर्थात् $x - y = 1$ और $x - y = -2$: फिर से असंगत। सारणिक का शून्य
होना यह घोषित करता है कि \emph{कुछ} अपभ्रष्ट हो रहा है, पर यह कभी नहीं बताता
कि क्या --- रिक्त या अनंत, यह दाएँ पक्ष को देखकर ही तय करना पड़ता है। यह भी
ध्यान दीजिए कि सूत्र स्वयं अपनी सीमाओं का संकेत दे देते हैं: $m \to 1^{-}$
होने पर $x = \frac{1 - 2m}{1 - m^2} \to -\infty$; अर्थात् जैसे-जैसे दोनों
रेखाएँ समांतर होती जाती हैं, हल-बिंदु भागता चला जाता है।
\end{example}

\begin{method}[निकायों पर गाउसीय विलोपन]\label{met:b1:det:gauss}
संवर्धित आव्यूह $(A \mid B)$ का पंक्ति-लघुकरण करके सोपानिक रूप तक
लाइए।\index{गाउसीय विलोपन}
\begin{enumerate}
  \item यदि अंतिम स्तंभ में कोई धुरी आ जाए ($0 = 1$ वाली पंक्ति): तो कोई हल
  नहीं।
  \item अन्यथा अज्ञात \emph{धुरी-अज्ञातों} और \emph{मुक्त अज्ञातों}
  (प्राचलों) में बँट जाते हैं; और पश्च-प्रतिस्थापन पहलों को दूसरों में
  व्यक्त कर देता है: \omterm{def:b1:logic:sets}{हल-समुच्चय} एक एफ़ीन \omterm{def:b1:vspaces:subspace}{उपसमष्टि} है जिसकी विमा $=$ मुक्त
  अज्ञातों की संख्या।
\end{enumerate}
क्रामर के सूत्र सिद्धांत और छोटे निकायों के लिए हैं; व्यावहारिक कलनविधि
विलोपन ही है।
\end{method}

\begin{example}[प्राचल के साथ एक विवेचन]\label{ex:b1:det:parameter}
$m \in \R$ के लिए विचार कीजिए
\[
\begin{cases}
x + y + mz = 1\\
x + my + z = 1\\
mx + y + z = 1 .
\end{cases}
\]
आव्यूह का सारणिक $-(m+2)(m-1)^2$ है (\cref{exo:b1:det:7} में सभी स्तंभ पहले
में जोड़कर संगणित)। $m \neq 1, -2$ के लिए: अद्वितीय हल $x = y = z =
\frac{1}{m+2}$ (सममिति से)। $m = 1$ के लिए: एक ही समीकरण तीन बार दोहराया
हुआ, अतः हलों का एक समतल। $m = -2$ के लिए: तीनों समीकरण जोड़ने पर $0 = 3$
मिलता है, अतः कोई हल नहीं।
\end{example}

\begin{remark}[सामान्य भूलें]
\label{rem:b1:det:pitfalls}
\emph{सारणिक आव्यूह में \omterm{def:b1:linmaps:def}{रैखिक} नहीं है}: $\det(A + B) \neq \det A + \det B$
(पहले से ही $\det(I_2 + I_2) = 4 \neq 2$); वह हर \emph{स्तंभ} में अलग-अलग
\omterm{def:b1:linmaps:def}{रैखिक} है, जो बिलकुल अलग बात है। \emph{मापन}: $\det(\lambda A) =
\lambda^n\det A$, $\lambda\det A$ नहीं --- $n$ स्तंभों में से हर एक मापित
होता है। \emph{सभी \omterm{met:b1:matrices:gauss}{पंक्ति संक्रियाएँ} मुफ़्त नहीं हैं}: $L_i \leftarrow L_i +
\lambda L_j$ सारणिक सुरक्षित रखता है, पर अदला-बदली चिह्न बदल देती है और $L_i
\leftarrow \lambda L_i$ उसे $\lambda$ से गुणा कर देता है --- यहाँ के
हिसाब-किताब की चूकें विलोपन-आधारित संगणनाओं में ग़लत चिह्नों का शास्त्रीय
स्रोत हैं। \emph{शून्य सारणिक आरंभ है, अंत नहीं}: वह कहता है ``कोटि $< n$'',
पर यह नहीं कि कौन-सी कोटि; उसका स्थान केवल आगे का काम बताता है (सोपानिक रूप,
अथवा \cref{exo:b1:det:12} के उपसारणिक) --- \cref{ex:b1:det:parameter} में $m
= 1$ बनाम $m = -2$ वाली स्थिति देखिए। \emph{क्रामर को व्युत्क्रमणीयता
चाहिए}: जब $\det A = 0$, तब सूत्र $x_j = \det A_j/\det A$ निरर्थक हैं, और
फिर भी निकाय के (अनंत कितने) हल भली-भाँति हो सकते हैं। \emph{सारणिक केवल
वर्ग आव्यूहों के होते हैं}: किसी आयताकार निकाय के लिए विलोपन ही एकमात्र
औज़ार है।
\end{remark}

\begin{remark}[सारणिक कहाँ जाते हैं]
\label{rem:b1:det:whereused}
इस अदिश की आगे तीन ज़िंदगियाँ हैं। \emph{ज्यामितीय}: $\abs{\det}$ संबद्ध
\omterm{def:b1:logic:map}{प्रतिचित्रण} का क्षेत्रफल- अथवा आयतन-मापन गुणक है --- जिसे समतल के लिए
\cref{ch:b1:euclid} में यथार्थ किया गया है, और चर-परिवर्तन के याकोबी के रूप
में स्नातक वर्ष~2 के खंड के बहु-समाकलों में। \emph{बीजगणितीय}: $\det(A -
\lambda I)$, अर्थात् \omterm{def:b1:diffeq:linear2}{अभिलक्षणिक बहुपद}, स्नातक वर्ष~2 में अभिलक्षणिक-मान
सिद्धांत खोल देता है --- और \cref{ch:b1:matrices} की सप्ताहांत समस्या की
सर्वसमिका $A^2 - (\operatorname{tr} A)A + (\det A)I = 0$ उसकी पहली छाया है।
\emph{विश्लेषणात्मक}: विशेष आव्यूहों (वांडरमोंड, कोशी, ग्राम) के सारणिक तय
करते हैं कि अंतर्वेशन, विघटन और \omterm{def:b1:linmaps:projection}{प्रक्षेप} की समस्याएँ कब सुप्रस्तुत हैं; और
नीचे की सप्ताहांत समस्या पहले दोनों कुलों का पूरा मूल्यांकन कर देती है।
\end{remark}

\begin{remark}[पुस्तक 3 के भीतर के परिदृश्य]
\label{rem:b1:det:perspectives}
यह अध्याय इस खंड की रैखिक-बीजगणित वाली रीढ़ बंद कर देता है, और शेष दो अध्याय
लाभांश भुनाते हैं। \cref{ch:b1:euclid} में: ग्राम आव्यूह $\bigl(\langle v_i,
v_j\rangle\bigr)$ किसी सारणिक से स्वतंत्रता की जाँच कर लेता है
(\cref{exo:b1:euclid:11}), और समतल के समदूरक अपने सारणिक के चिह्न के अनुसार
घूर्णनों और परावर्तनों में बँट जाते हैं --- वहाँ की सप्ताहांत समस्या का
वर्गीकरण इसी पर चलता है। \cref{ch:b1:multivar} में: मोंज राशि $rt - s^2$
द्वितीय अवकलजों के सममित आव्यूह का सारणिक है, और न्यूनतम वर्गों के
प्रसामान्य समीकरण एक ऐसा क्रामर निकाय हैं जिसका आव्यूह ग्राम (अतः आघूर्ण)
आव्यूह है --- और वह ठीक यहीं स्थापित वांडरमोंड-स्वाद वाली कसौटियों से
व्युत्क्रमणीय होता है। जब वे अध्याय ``व्युत्क्रमणीय'' या ``धनात्मक'' कहते
हैं, तो उसकी रसीदें इसी अध्याय में हैं।
\end{remark}

\section{अभ्यास}

\begin{exercise}[$\star$]\label{exo:b1:det:1}
संगणित कीजिए:
\[
\begin{vmatrix} 3 & 1\\ 5 & 2 \end{vmatrix},
\qquad
\begin{vmatrix} 1 & 2 & 3\\ 4 & 5 & 6\\ 7 & 8 & 9\end{vmatrix},
\qquad
\begin{vmatrix} 1 & 1 & 1\\ 1 & 2 & 4\\ 1 & 3 & 9\end{vmatrix}.
\]
\end{exercise}

\begin{exercise}[$\star$]\label{exo:b1:det:2}
किन $\lambda \in \R$ के लिए कुल $\bigl((1, 1, \lambda), (1, \lambda, 1),
(\lambda, 1, 1)\bigr)$ $\R^3$ का \omterm{def:b1:vspaces:free}{आधार} है?
\end{exercise}

\begin{exercise}[$\star$]\label{exo:b1:det:3}
क्रामर के नियम से हल कीजिए:
\[
\begin{cases}
2x + y = 5\\
3x - 2y = 4 ,
\end{cases}
\qquad\text{फिर}\qquad
\begin{cases}
x + y + z = 6\\
x - y + z = 2\\
2x + y - z = 1 .
\end{cases}
\]
\end{exercise}

\begin{exercise}[$\star$]\label{exo:b1:det:4}
\omterm{met:b1:det:gauss}{गाउसीय विलोपन} से हल कीजिए, और \omterm{def:b1:logic:sets}{हल-समुच्चय} का वर्णन कीजिए:
\[
\begin{cases}
x + 2y - z + t = 1\\
2x + 4y + z - t = 5\\
x + 2y + 2z - 2t = 4 .
\end{cases}
\]
\end{exercise}

\begin{exercise}[$\star\star$]\label{exo:b1:det:5}
\cref{ex:b1:det:vandermonde} का वांडरमोंड सूत्र $n$ पर आगमन से सिद्ध कीजिए,
और स्तंभ संक्रियाएँ $C_k \leftarrow C_k - x_1 C_{k-1}$ $k = n$ से $k = 2$ तक
नीचे की ओर कीजिए।
\end{exercise}

\begin{exercise}[$\star\star$]\label{exo:b1:det:6}
(त्रिविकर्ण) मान लीजिए $D_n$ वह $n \times n$ सारणिक है जिसके विकर्ण पर $2$,
दोनों निकटवर्ती विकर्णों पर $1$, और अन्यत्र $0$ है। पहली पंक्ति के अनुदिश
प्रसार करके $D_n = 2D_{n-1} - D_{n-2}$ सिद्ध कीजिए और $D_n$ संगणित कीजिए
($D_1 = 2$, $D_2 = 3$)।
\end{exercise}

\begin{exercise}[$\star\star$]\label{exo:b1:det:7}
\cref{ex:b1:det:parameter} को पूरा कीजिए: संक्रिया $C_1 \leftarrow C_1 + C_2
+ C_3$ से सारणिक $\begin{vmatrix} 1 & 1 & m\\ 1 & m & 1\\ m & 1 &
1\end{vmatrix}$ संगणित कीजिए, और निकाय का पूरा विवेचन कीजिए।
\end{exercise}

\begin{exercise}[$\star\star$]\label{exo:b1:det:8}
मान लीजिए $A \in \mathcal{M}_n(\R)$, जिसकी प्रविष्टियाँ \emph{पूर्णांक} हैं।
सिद्ध कीजिए कि $A$ का पूर्णांक प्रविष्टियों वाला प्रतिलोम है तभी जब $\det A
= \pm 1$। \emph{(सीधी दिशा के लिए सारणिक लीजिए; विलोम के लिए यह मान लीजिए
--- अथवा $n \leq 3$ के लिए सहगुणनखंडों से सिद्ध कीजिए --- कि $A^{-1} =
\frac{1}{\det A}\,\operatorname{Com}(A)^{\mathsf T}$, जहाँ सहगुणनखंड आव्यूह
पूर्णांक वाला है।)}
\end{exercise}

\begin{exercise}[$\star\star\star$]\label{exo:b1:det:9}
आव्यूह $aI + bJ$ (\cref{exo:b1:matrices:9}) का $n \times n$ सारणिक संगणित
कीजिए, अर्थात् जिसके विकर्ण पर $a + b$ और अन्यत्र $b$ हो। \emph{(सभी स्तंभ
पहले में जोड़िए, गुणनखंडन कीजिए, फिर साफ़ कीजिए।)} व्युत्क्रमणीयता की शर्त
$a \neq 0$, $a + nb \neq 0$ फिर से प्राप्त कीजिए।
\end{exercise}

\begin{exercise}[$\star\star\star$]\label{exo:b1:det:10}
मान लीजिए $A, B \in \mathcal{M}_n(\R)$। सिद्ध कीजिए कि
\[
\det\begin{pmatrix} A & B\\ B & A \end{pmatrix}
= \det(A + B)\,\det(A - B),
\]
इसके लिए खंड-स्तंभ और खंड-पंक्ति संक्रियाओं का प्रयोग कीजिए (खंड-रूप में
$C_1 \leftarrow C_1 + C_2$, फिर $L_2 \leftarrow L_2 - L_1$), और स्वाभाविक
खंड-त्रिभुजाकार नियम $\det\begin{pmatrix} M & N\\ 0 & P\end{pmatrix} = \det
M \det P$ मान लीजिए --- जो \cref{ex:b1:det:blocktriangular} में $2 \times 2$
खंडों के लिए सिद्ध है।
\end{exercise}

\begin{exercise}[$\star\star$]\label{exo:b1:det:11}
(कोटि $3$ का चक्रीय आव्यूह) मान लीजिए $a, b, c \in \C$ और
\[
\Delta = \begin{vmatrix}
a & b & c\\
c & a & b\\
b & c & a
\end{vmatrix}.
\]
सिद्ध कीजिए कि $\Delta = (a + b + c)(a^2 + b^2 + c^2 - ab - bc - ca)$, और $j
= \eu^{2\iu\pi/3}$ का प्रयोग करके $\C$ पर पूरा गुणनखंडन कीजिए:
\[
\Delta = (a + b + c)(a + jb + j^2c)(a + j^2b + jc) .
\]
\emph{($C_1 \leftarrow C_1 + C_2 + C_3$ से आरंभ कीजिए; सम्मिश्र रूप के लिए
ध्यान दीजिए कि स्तंभ $(1, j, j^2)^{\mathsf T}$ लगभग किसी अभिलक्षणिक सदिश की
तरह व्यवहार करता है।)}
\end{exercise}

\begin{exercise}[$\star\star\star$]\label{exo:b1:det:12}
(कोटि और उपसारणिक) मान लीजिए $A \in \mathcal{M}_{n,p}(K)$। सिद्ध कीजिए कि
$\operatorname{rk} A$ $A$ के किसी व्युत्क्रमणीय $r \times r$ उपआव्यूह के
अधिकतम आकार $r$ के बराबर है (उपआव्यूह $r$ चुनी हुई पंक्तियों और $r$ चुने हुए
स्तंभों के कटान-बिंदुओं की प्रविष्टियाँ रखता है)। \emph{(यदि
$\operatorname{rk} A = r$, तो $r$ \omterm{def:b1:vspaces:free}{स्वतंत्र} स्तंभ चुनिए, फिर प्राप्त $n
\times r$ खंड की $r$ \omterm{def:b1:vspaces:free}{स्वतंत्र} पंक्तियाँ; विलोमतः कोई व्युत्क्रमणीय उपआव्यूह
$A$ के तदनुरूप स्तंभों को \omterm{def:b1:vspaces:free}{स्वतंत्र} होने पर बाध्य कर देता है।)}
\end{exercise}

\section{समस्या: कोशी का दुहरा एकांतरक}

\begin{problem}\label{pb:b1:det:1}
इस अध्याय के अनुप्रयोगों पर दो सारणिक राज करते हैं: \omterm{ex:b1:det:vandermonde}{वांडरमोंड सारणिक}, जिसका
मूल्यांकन \cref{exo:b1:det:5} में हुआ, और \emph{\omterm{pb:b1:det:1}{कोशी सारणिक}}
$\det\bigl(\frac{1}{a_i + b_j}\bigr)$, जिसका मूल्यांकन यहाँ है। उन्हीं के
इर्द-गिर्द यह समस्या एकांतरक का औज़ार-बक्सा इकट्ठा करती है: बहुपदीय
स्तंभ-चालें, क्रामर से अंतर्वेशन, \omterm{pb:b1:det:1}{हिल्बर्ट आव्यूह}, किसी त्रिघातीय का
विविक्तकर, और एकांतरक बहुपदों की विधि। सर्वत्र $V(x_1, \dots, x_n) =
\prod_{i < j}(x_j - x_i)$ वांडरमोंड मान को दर्शाता है।\index{कोशी
सारणिक}\index{हिल्बर्ट आव्यूह} \index{विविक्तकर}

\medskip \noindent\textbf{भाग I --- वांडरमोंड का औज़ार-बक्सा।}
\begin{enumerate}
  \item $V(1, 2, 3, 4)$ संगणित कीजिए, और स्मरण कीजिए कि $n$ युग्मशः भिन्न
  गाँठों पर अंतर्वेशन क्रामर निकाय क्यों है।
  \item (बहुपदीय एकांतरक) मान लीजिए $P_0, \dots, P_{n-1}$ \emph{\omterm{def:b1:poly:def}{इकाई-अग्र}}
  हैं और $\deg P_k = k$। सिद्ध कीजिए
        \[
        \det\bigl(P_{j-1}(x_i)\bigr)_{1 \leq i, j \leq n}
        = V(x_1, \dots, x_n) :
        \]
        अर्थात् स्तंभ संक्रियाएँ हर घात-स्तंभ को किसी भी \omterm{def:b1:poly:def}{इकाई-अग्र} सीढ़ी से
        मुफ़्त में बदल देती हैं।
  \item प्रश्न 2 को द्विपद बहुपदों $B_k = \frac{X(X-1)\cdots(X-k+1)}{k!}$ पर
  लगाइए: सिद्ध कीजिए कि \emph{पूर्णांकों} $m_1 < m_2 < \dots < m_n$ के लिए,
        \[
        \frac{V(m_1, \dots, m_n)}{0!\,1!\,2!\cdots(n-1)!}
        \in \N :
        \]
        अर्थात् $n$ पूर्णांकों के सभी युग्मशः अंतरों का गुणनफल अधिक्रमगुणित
        $0!\,1!\cdots(n-1)!$ से विभाज्य होता है।
  \item $\det\bigl(x_i^{\,j}\bigr)_{1 \leq i, j \leq n} = x_1 x_2 \cdots
  x_n\, V(x_1, \dots, x_n)$ सिद्ध कीजिए (अब घातें $1$ से आरंभ होती हैं)।
  \item (आघूर्ण आव्यूह) मान लीजिए $S = \bigl(p_{i+j-2}\bigr)_{1 \leq i, j
  \leq n}$, जहाँ $p_k = x_1^k + \dots + x_n^k$। सिद्ध कीजिए कि आव्यूह $W =
  (x_i^{\,j-1})_{ij}$ के लिए $S = W^{\mathsf T} W$, और निकालिए
        \[
        \det S = V(x_1, \dots, x_n)^2 ,
        \]
        और निष्कर्ष निकालिए: $n$ \emph{वास्तविक} संख्याएँ युग्मशः भिन्न हैं
        तभी जब उनका आघूर्ण आव्यूह व्युत्क्रमणीय हो, और $\det S \geq 0$ सदा।
\end{enumerate}

\medskip \noindent\textbf{भाग II --- अंतर्वेशन, फिर से।} गाँठें $x_1 < \dots
< x_n$, मान $y_1, \dots, y_n$।
\begin{enumerate}[resume]
  \item शर्तों ``$P = c_0 + c_1X + \dots + c_{n-1}X^{n-1}$ अंतर्वेशन करता
  है'' को $c_k$ में आव्यूह $W$ वाले \omterm{def:b1:det:system}{रैखिक निकाय} के रूप में लिखिए, और $\det W
  = V \neq 0$ से अंतर्वेशक का अस्तित्व तथा अद्वितीयता फिर से प्राप्त कीजिए
  (पहले की दोनों उपपत्तियों \cref{thm:b1:poly:lagrange} और
  \cref{ex:b1:linmaps:interpolation} से तुलना कीजिए)।
  \item क्रामर के नियम से और संबंधित सारणिक के अंतिम स्तंभ के अनुदिश
  \omterm{thm:b1:det:cofactor}{सहगुणनखंड-प्रसार} से सिद्ध कीजिए कि अंतर्वेशक का अग्र गुणांक है
        \[
        c_{n-1} = \sum_{i=1}^{n}
        \frac{y_i}{\prod_{j \neq i}(x_i - x_j)} .
        \]
  \item (संगामी वांडरमोंड) संगणित कीजिए
        \[
        \begin{vmatrix}
        1 & x_1 & x_1^2\\
        0 & 1 & 2x_1\\
        1 & x_2 & x_2^2
        \end{vmatrix}
        = (x_2 - x_1)^2 ,
        \]
        और व्याख्या कीजिए: $x_1 \neq x_2$ होने पर आँकड़े $\bigl(P(x_1),
        P'(x_1), P(x_2)\bigr)$ एक अद्वितीय $P \in \R_2[X]$ तय कर देते हैं
        (एरमीट अंतर्वेशन)।
  \item $P(0) = 1$, $P'(0) = 0$, $P(1) = 2$ वाला अद्वितीय $P \in \R_2[X]$
  खोजिए, और अपना उत्तर प्रश्न 8 के सामने जाँचिए।
\end{enumerate}

\medskip \noindent\textbf{भाग III --- \omterm{pb:b1:det:1}{कोशी सारणिक}।} मान लीजिए $a_1, \dots,
a_n$ और $b_1, \dots, b_n$ ऐसे अदिश हैं कि सभी $i, j$ के लिए $a_i + b_j \neq
0$, और
\[
C_n = \det\Bigl(\frac{1}{a_i + b_j}\Bigr)_{1 \leq i, j \leq n} .
\]
\begin{enumerate}[resume]
  \item $C_2$ हाथ से संगणित कीजिए और उसे ``अंतरों के गुणनफल बटा योगों के
  गुणनफल'' के रूप में रखिए।
  \item $n \geq 2$ के लिए $L_i \leftarrow L_i - L_n$ कीजिए ($i < n$) और
  पंक्तियों तथा स्तंभों का गुणनखंडन करके सिद्ध कीजिए
        \[
        C_n = \frac{\prod_{i<n}(a_n - a_i)}{\prod_{j}(a_n +
        b_j)}\;\det M,
        \]
        जहाँ $M$ पंक्तियों $i < n$ पर कोशी आव्यूह से मेल खाता है और उसकी
        अंतिम पंक्ति $(1, 1, \dots, 1)$ है।
  \item $M$ पर $C_j \leftarrow C_j - C_n$ कीजिए ($j < n$), फिर से गुणनखंडन
  कीजिए, और आगमन से \emph{कोशी का दुहरा एकांतरक} निष्कर्ष निकालिए:
        \[
        C_n = \frac{\prod_{1 \leq i < j \leq n}(a_j - a_i)(b_j -
        b_i)}{\prod_{i, j}(a_i + b_j)} .
        \]
  \item व्युत्क्रमणीयता की कसौटी निकालिए ($a_i$ युग्मशः भिन्न और $b_j$
  युग्मशः भिन्न)। \emph{\omterm{pb:b1:det:1}{हिल्बर्ट आव्यूह}} $H_n = \bigl(\frac{1}{i + j -
  1}\bigr)$ के लिए: सूत्र से $\det H_2$ और $\det H_3$ संगणित कीजिए, और
  सत्यापित कीजिए कि $H_2^{-1}$ की प्रविष्टियाँ पूर्णांक हैं।
  \item दिखाइए कि युग्मशः भिन्न $b_j$ और किसी भी दाएँ पक्ष के लिए निकाय
  $\sum_j \frac{c_j}{a_i + b_j} = y_i$ ($i = 1, \dots, n$) का अद्वितीय हल
  है, और इसे सरल ध्रुवों वाले आंशिक भिन्न-विघटनों के अस्तित्व तथा अद्वितीयता
  से जोड़िए (\cref{thm:b1:fractions:complex})।
\end{enumerate}

\medskip \noindent\textbf{भाग IV --- त्रिघातीय का विविक्तकर।} मान लीजिए
$\lambda_1, \lambda_2, \lambda_3$ $X^3 + pX + q$ के ($\C$ में) मूल हैं, और
$p_k = \lambda_1^k + \lambda_2^k + \lambda_3^k$।
\begin{enumerate}[resume]
  \item हर मूल पर $\lambda^3 = -p\lambda - q$ का और वियेत ($p_1 = 0$) का
  प्रयोग करके $p_2 = -2p$, $p_3 = -3q$ तथा $p_4 = 2p^2$ संगणित कीजिए।
  \item प्रश्न 5 के साथ ($\C$ पर, $\det S = V^2$ रखते हुए) संगणित कीजिए
        \[
        \operatorname{disc} = V(\lambda_1, \lambda_2,
        \lambda_3)^2 = \begin{vmatrix}
        3 & 0 & -2p\\
        0 & -2p & -3q\\
        -2p & -3q & 2p^2
        \end{vmatrix}
        = -4p^3 - 27q^2 .
        \]
  \item निष्कर्ष निकालिए: $X^3 + pX + q$ का कोई दोहराया मूल है तभी जब $4p^3
  + 27q^2 = 0$; $X^3 - 3X + 2 = (X - 1)^2(X + 2)$ पर जाँचिए।
  \item मान लीजिए $p, q$ वास्तविक हैं। सिद्ध कीजिए कि त्रिघातीय के तीन भिन्न
  वास्तविक मूल हैं तभी जब $\operatorname{disc} > 0$, और एक वास्तविक तथा दो
  अवास्तविक संयुग्मी मूल हैं तभी जब $\operatorname{disc} < 0$। \emph{(यदि
  $\lambda_3 = \conj{\lambda_2} \neq \lambda_2$ और $\lambda_1 \in \R$, तो
  दिखाइए कि $V$ विशुद्ध काल्पनिक है।)}
\end{enumerate}

\medskip \noindent\textbf{भाग V --- लाभांश, और एकांतरक विधि।}
\begin{enumerate}[resume]
  \item $0 < a_1 < a_2 < \dots < a_n$ के लिए दिखाइए $\det\bigl(\frac{1}{a_i
  + a_j}\bigr) > 0$।
  \item $(m_1, m_2, m_3) = (2, 4, 7)$ के लिए
  $\det\bigl(\binom{m_i}{j-1}\bigr)_{1 \leq i, j \leq 3}$ संगणित कीजिए, पहले
  प्रश्न 2--3 से, फिर सीधे प्रसार से।
  \item मान लीजिए $\lambda_1, \dots, \lambda_n$ युग्मशः भिन्न और अशून्य हैं।
  किसी व्युत्क्रमणीय वांडरमोंड आव्यूह का प्रयोग करके फिर से सिद्ध कीजिए कि
  गुणोत्तर अनुक्रम $\bigl((\lambda_i^{\,k})_{k \geq 0}\bigr)_{1 \leq i \leq
  n}$ अनुक्रमों की समष्टि का \omterm{def:b1:vspaces:free}{स्वतंत्र कुल} बनाते हैं।
  \item दुहरे एकांतरक से $\det\bigl(\frac{1}{i + j}\bigr)_{1 \leq i, j \leq
  3}$ संगणित कीजिए।
  \item (एकांतरक \omterm{def:b1:poly:def}{बहुपद}) $x_1, \dots, x_n$ में किसी \omterm{def:b1:poly:def}{बहुपद} $F$ को
  \emph{एकांतरक} कहिए जब किन्हीं दो चरों की अदला-बदली उसका चिह्न बदल दे।
  दिखाइए कि $x_i = x_j$ होते ही एकांतरक $F$ शून्य हो जाता है ($i \neq j$),
  और --- गुणनखंड प्रमेय से एक-एक चर लेकर --- निष्कर्ष निकालिए कि $F$
  $\prod_{i<j}(x_j - x_i)$ से विभाज्य है।
  \item प्रश्न 23 का प्रयोग करके वांडरमोंड सूत्र बिना किसी आगमन के फिर से
  सिद्ध कीजिए: सारणिक $\det(x_i^{\,j-1})$ कुल घात $\binom n2$ का एकांतरक
  \omterm{def:b1:poly:def}{बहुपद} है, अतः $\prod_{i<j}(x_j - x_i)$ का कोई \emph{अचर} गुणज; और एक एकपदी
  की तुलना करके वह अचर पहचानिए।
  \item संश्लेषण, चार वाक्यों में: सारणिक का कौन-सा अकेला गुण (कौन-सा
  अभिगृहीत) इस समस्या के सभी गुणनखंडन पैदा करता है; प्रश्न 5 की
  आघूर्ण-आव्यूह सर्वसमिका \emph{सम्मिश्र} भिन्नता के \omterm{def:b1:logic:statement}{कथन} को संगणनीय
  \emph{वास्तविक} चिह्न-परीक्षा में क्यों बदल देती है; यहाँ कौन-से दो
  शास्त्रीय आव्यूहों का पूरा मूल्यांकन हुआ और वे कौन-सी \omterm{def:b1:linmaps:def}{रैखिक} समस्याओं पर
  शासन करते हैं; और प्रश्न 23--24 की एकांतरक विधि एक ही झटके में यह क्यों
  समझा देती है कि $\prod_{i<j}(x_j - x_i)$ बार-बार क्यों आता रहता है। भाग
  III की प्रमेय का नाम बताइए।
\end{enumerate}
\end{problem}
